\subsection{Quantitative signature of projection entropy}
  \label{subsec:projection-entropy-bound}

  The structural analysis developed above identifies non-injectivity as a sufficient
  condition for the failure of Bell factorizability.
  We now propose a minimal quantitative characterization of this effect.

  Let
  \begin{equation}
    S_{\Pi} \equiv - \sum_{o \in O} \nu(o)\,\log \nu(o),
  \end{equation}
  denote the projection entropy associated with the mapping $\Pi : \Omega \to O$,
  where $\nu(o) = \mu(\Pi^{-1}(o))$ is the induced measure on observable outcomes.

  The quantity $S_{\Pi}$ measures the effective multiplicity of underlying configurations
  associated with each observable state.
  It therefore provides a natural scalar measure of non-injectivity.

  We propose that the strength of Bell--CHSH violations is controlled by $S_{\Pi}$.
  In particular, the CHSH parameter
  \begin{equation}
    S \equiv |E_{00} + E_{01} + E_{10} - E_{11}|
  \end{equation}
  is expected to satisfy a structural bound of the form
  \begin{equation}
    S \;\leq\; 2 + \Delta(S_{\Pi}),
  \end{equation}
  where $\Delta(S_{\Pi})$ is a monotonically increasing function satisfying
  \begin{equation}
    \Delta(0) = 0.
  \end{equation}

  This expresses the fact that Bell violations arise only in regimes where the
  projection entropy is non-zero.
  In the limit of an injective mapping ($S_{\Pi} = 0$), standard classical bounds are recovered.

  Conversely, increasing projection entropy corresponds to an increasing overlap
  between equivalence classes of underlying configurations, enhancing the degree
  of non-factorizability of observable probabilities.

  The explicit form of $\Delta(S_{\Pi})$ depends on additional structural constraints
  on the underlying configuration space and the projection.
  In particular, physically realizable theories may impose bounds on admissible
  correlations that restrict $S$ below its algebraic maximum.

  This suggests that quantum correlations may correspond to an intermediate regime
  in which $S_{\Pi}$ is non-zero but constrained, leading to bounded violations
  consistent with the Tsirelson limit.

  The relation between projection entropy and Bell-type violations thus provides
  a potential quantitative bridge between structural non-injectivity and observable
  correlation strength, and defines a testable signature of the present framework.
